%%%  کلاس AUTthesis، نسخه آبان 1397
%%%   دانشگاه صنعتی امیرکبیر                 http://www.aut.ac.ir
%%%  تالار گفتگوی پارسی‌لاتک،       http://forum.parsilatex.com
%%%   آپدیت شده در آبان 95
%%%   پشتیبانی و راهنمایی          badali_farhad@yahoo.com
%%%
%%%   بازبینی و اصلاح شده در آبان ماه 1397
%%%  Tested via TeXstudio in TeXlive 2014-2018.
%%%

%-----------------------------------------------------------------------------------------------------
%        روش اجرا.: 2 بار F1 ، 2 بار  F11(به منظور تولید مراجع) ، دوبار Ctrl+Alt+I (به منظور تولید نمایه) و دو بار F1 -------> مشاهده Pdf
%%%%%%%%%%%%%%%%%%%%%%%%%%%%%%%%%%%%%%%%%%%%%%%%%%%%%%
%   TeXstudio as your IDE
%%  برای compile در TeXstudio تنها کافی است منوی Options->Configure TeXstudio را زده و در پنجره Configure TeXstudio در بخش Build گزینه Default Compiler را به XeLaTeX تغییر دهید. سند شما به راحتی compile خواهد شد.
%   F1 & F5 : Build & view
%   F6      : Compile
%   F7      : Viewira
%   --------------
%%%%%%%%%%%%%%%%%%%%%%%%%%%%%%%%%%%%%%%%%%%%%%%%%%%%%%
%        اگر قصد نوشتن رساله دکتری را دارید، در خط زیر به جای msc،
%      کلمه phd را قرار دهید. کلیه تنظیمات لازم، به طور خودکار، اعمال می‌شود.
%%% !TEX TS-program = XeLaTeX

\documentclass[oneside,msc,12pt]{AUTthesis}
\usepackage{pgf}

\usepackage{pythonhighlight}

\usepackage{listings}
\usepackage{color}

\definecolor{dkgreen}{rgb}{0,0.6,0}
\definecolor{gray}{rgb}{0.5,0.5,0.5}
\definecolor{mauve}{rgb}{0.58,0,0.82}

\lstset{frame=tb,
  language=Python,
  aboveskip=3mm,
  belowskip=3mm,
  showstringspaces=false,
  columns=flexible,
  basicstyle={\small\ttfamily},
  numbers=none,
  numberstyle=\tiny\color{gray},
  keywordstyle=\color{blue},
  commentstyle=\color{dkgreen},
  stringstyle=\color{mauve},
  breaklines=true,
  breakatwhitespace=true,
  tabsize=3
}


%       فایل commands.tex را حتماً به دقت مطالعه کنید؛ چون دستورات مربوط به فراخوانی بسته زی‌پرشین 
%       و دیگر بسته‌ها و ... در این فایل قرار دارد و بهتر است که با نحوه استفاده از آنها آشنا شوید. توجه شود برای نسخه نهایی پایان‌نامه حتماً hyperref را 
%        غیرفعال کنید.


% در این فایل، دستورها و تنظیمات مورد نیاز، آورده شده است.
%-------------------------------------------------------------------------------------------------------------------
% در ورژن جدید زی‌پرشین برای تایپ متن‌های ریاضی، این سه بسته، حتماً باید فراخوانی شود.
\usepackage{amsthm,amssymb,amsmath,amsfonts}
% بسته‌ای برای تنطیم حاشیه‌های بالا، پایین، چپ و راست صفحه
\usepackage[top=30mm, bottom=30mm, left=25mm, right=30mm]{geometry}
% بسته‌‌ای برای ظاهر شدن شکل‌ها و تصاویر متن
\usepackage{graphicx}
\usepackage{color}
%بسته‌ای برای تنظیم فاصله عمودی خط‌های متن
\usepackage{setspace}
\usepackage{titletoc}
\usepackage{tocloft}
%با فعال کردن بسته زیر فوت‌نوت‌ها در هر صفحه ریست می‌شوند. حالت پیش‌فرض آن ریست شدن در هر فصل می‌باشد.
%\usepackage[perpage]{footmisc}
\usepackage{enumitem}
%\usepackage{titlesec}
% بسته‌ و دستوراتی برای ایجاد لینک‌های رنگی با امکان جهش
\usepackage[pagebackref=false,colorlinks,linkcolor=blue,citecolor=red]{hyperref}
\usepackage[nameinlink]{cleveref}%capitalize,,noabbrev
 \AtBeginDocument{%
    \crefname{equation}{برابری}{equations}%
    \crefname{chapter}{فصل}{chapters}%
    \crefname{section}{بخش}{sections}%
    \crefname{appendix}{پیوست}{appendices}%
    \crefname{enumi}{مورد}{items}%
    \crefname{footnote}{زیرنویس}{footnotes}%
    \crefname{figure}{شکل}{figures}%
    \crefname{table}{جدول}{tables}%
    \crefname{theorem}{قضیه}{theorems}%
    \crefname{lemma}{لم}{lemmas}%
    \crefname{corollary}{نتیجه}{corollaries}%
    \crefname{proposition}{گزاره}{propositions}%
    \crefname{definition}{تعریف}{definitions}%
    \crefname{result}{نتیجه}{results}%
    \crefname{example}{مثال}{examples}%
    \crefname{remark}{نکته}{remarks}%
    \crefname{note}{یادداشت}{notes}%
}
% چنانچه قصد پرینت گرفتن نوشته خود را دارید، خط بالا را غیرفعال و  از دستور زیر استفاده کنید چون در صورت استفاده از دستور زیر‌‌، 
% لینک‌ها به رنگ سیاه ظاهر خواهند شد که برای پرینت گرفتن، مناسب‌تر است
%\usepackage[pagebackref=false]{hyperref}
% بسته‌ لازم برای تنظیم سربرگ‌ها
\usepackage{fancyhdr}
% بسته‌ای برای ظاهر شدن «مراجع»  در فهرست مطالب
\usepackage[nottoc]{tocbibind}
% دستورات مربوط به ایجاد نمایه
\usepackage{makeidx,multicol}
\setlength{\columnsep}{1.5cm}

%%%%%%%%%%%%%%%%%%%%%%%%%%
\usepackage{verbatim}
\makeindex
\usepackage{sectsty}
% فراخوانی بسته زی‌پرشین و تعریف قلم فارسی و انگلیسی
\usepackage{xepersian}%[extrafootnotefeatures]
\SepMark{-}
%حتماً از تک لایو 2014 استفاده کنید.
\settextfont[Scale=1.2]{B Nazanin}
\setlatintextfont{Times New Roman}
\setsansfont{Trebuchet MS} 
\setmonofont{mononoki Nerd Font Mono}
\renewcommand{\labelitemi}{$\bullet$}
%%%%%%%%%%%%%%%%%%%%%%%%%%
% چنانچه می‌خواهید اعداد در فرمول‌ها، انگلیسی باشد، خط زیر را غیرفعال کنید.
%در غیر اینصورت حتماً فونت PGaramond را نصب کنید.
%\setdigitfont[Scale=1.1]{PGaramond}%%Yas
%%%%%%%%%%%%%%%%%%%%%%%%%%
% تعریف قلم‌های فارسی اضافی برای استفاده در بعضی از قسمت‌های متن
\defpersianfont\nastaliq[Scale=2]{IranNastaliq}
\defpersianfont\chapternumber[Scale=3]{B Nazanin}
%\chapterfont{\centering}%
%%%%%%%%%%%%%%%%%%%%%%%%%%
% دستوری برای تغییر نام کلمه «اثبات» به «برهان»
\renewcommand\proofname{\textbf{برهان}}

% دستوری برای تغییر نام کلمه «کتاب‌نامه» به «منابع و مراجع«
\renewcommand{\bibname}{منابع و مراجع}


% Headings for every page of ToC, LoF and Lot
\setlength{\cftbeforetoctitleskip}{-1.2em}
\setlength{\cftbeforelottitleskip}{-1.2em}
\setlength{\cftbeforeloftitleskip}{-1.2em}
\setlength{\cftaftertoctitleskip}{-1em}
\setlength{\cftafterlottitleskip}{-1em}
\setlength{\cftafterloftitleskip}{-1em}
%%\makeatletter
%%%%\renewcommand{\l@chapter}{\@dottedtocline{1}{1em\bfseries}{1em}}
%%%%\renewcommand{\l@section}{\@dottedtocline{2}{2em}{2em}}
%%%%\renewcommand{\l@subsection}{\@dottedtocline{3}{3em}{3em}}
%%%%\renewcommand{\l@subsubsection}{\@dottedtocline{4}{4em}{4em}}
%%%%\makeatother


\newcommand\tocheading{\par عنوان\hfill صفحه \par}
\newcommand\lofheading{\hspace*{.5cm}\figurename\hfill صفحه \par}
\newcommand\lotheading{\hspace*{.5cm}\tablename\hfill صفحه \par}

\renewcommand{\cftchapleader}{\cftdotfill{\cftdotsep}}
\renewcommand{\cfttoctitlefont}{\hspace*{\fill}\LARGE\bfseries}%\Large
\renewcommand{\cftaftertoctitle}{\hspace*{\fill}}
\renewcommand{\cftlottitlefont}{\hspace*{\fill}\LARGE\bfseries}%\Large
\renewcommand{\cftafterlottitle}{\hspace*{\fill}}
\renewcommand{\cftloftitlefont}{\hspace*{\fill}\LARGE\bfseries}
\renewcommand{\cftafterloftitle}{\hspace*{\fill}}

%%%%%%%%%%%%%%%%%%%%%%%%%%
% تعریف و نحوه ظاهر شدن عنوان قضیه‌ها، تعریف‌ها، مثال‌ها و ...
%برای شماره گذاری سه تایی قضیه ها
\theoremstyle{definition}
\newtheorem{definition}{تعریف}[section]
\newtheorem{remark}[definition]{نکته}
\newtheorem{note}[definition]{یادداشت}
\newtheorem{example}[definition]{نمونه}
\newtheorem{question}[definition]{سوال}
\newtheorem{remember}[definition]{یاداوری}
\theoremstyle{theorem}
\newtheorem{theorem}[definition]{قضیه}
\newtheorem{lemma}[definition]{لم}
\newtheorem{proposition}[definition]{گزاره}
\newtheorem{corollary}[definition]{نتیجه}
%%%%%%%%%%%%%%%%%%%%%%%%
%%%%%%%%%%%%%%%%%%%
%%% برای شماره گذاری چهارتایی قضیه ها و ...
%%\newtheorem{definition1}[subsubsection]{تعریف}
%%\newtheorem{theorem1}[subsubsection]{قضیه}
%%\newtheorem{lemma1}[subsubsection]{لم}
%%\newtheorem{proposition1}[subsubsection]{گزاره}
%%\newtheorem{corollary1}[subsubsection]{نتیجه}
%%\newtheorem{remark1}[subsubsection]{نکته}
%%\newtheorem{example1}[subsubsection]{مثال}
%%\newtheorem{question1}[subsubsection]{سوال}

%%%%%%%%%%%%%%%%%%%%%%%%%%%%

% دستورهایی برای سفارشی کردن صفحات اول فصل‌ها
\makeatletter
\newcommand\mycustomraggedright{%
 \if@RTL\raggedleft%
 \else\raggedright%
 \fi}
\def\@makechapterhead#1{%
\thispagestyle{style1}
\vspace*{20\p@}%
{\parindent \z@ \mycustomraggedright
\ifnum \c@secnumdepth >\m@ne
\if@mainmatter

\bfseries{\Huge \@chapapp}\small\space {\chapternumber\thechapter}
\par\nobreak
\vskip 0\p@
\fi
\fi
\interlinepenalty\@M 
\Huge \bfseries #1\par\nobreak
\vskip 120\p@

}

%\thispagestyle{empty}
\newpage}
\bidi@patchcmd{\@makechapterhead}{\thechapter}{\tartibi{chapter}}{}{}
\bidi@patchcmd{\chaptermark}{\thechapter}{\tartibi{chapter}}{}{}
\makeatother

\pagestyle{fancy}
\renewcommand{\chaptermark}[1]{\markboth{\chaptername~\tartibi{chapter}: #1}{}}

\fancypagestyle{style1}{
\fancyhf{} 
\fancyfoot[c]{\thepage}
\fancyhead[R]{\leftmark}%
\renewcommand{\headrulewidth}{1.2pt}
}


\fancypagestyle{style2}{
\fancyhf{}
\fancyhead[R]{چکیده}
\fancyfoot[C]{\thepage{}}
\renewcommand{\headrulewidth}{1.2pt}
}

\fancypagestyle{style3}{%
  \fancyhf{}%
  \fancyhead[R]{فهرست نمادها}
  \fancyfoot[C]{\thepage}%
  \renewcommand{\headrulewidth}{1.2pt}%
}

\fancypagestyle{style4}{%
  \fancyhf{}%
  \fancyhead[R]{فهرست جداول}
  \fancyfoot[C]{\thepage}%
  \renewcommand{\headrulewidth}{1.2pt}%
}

\fancypagestyle{style5}{%
  \fancyhf{}%
  \fancyhead[R]{فهرست اشکال}
  \fancyfoot[C]{\thepage}%
  \renewcommand{\headrulewidth}{1.2pt}%
}

\fancypagestyle{style6}{%
  \fancyhf{}%
  \fancyhead[R]{فهرست مطالب}
  \fancyfoot[C]{\thepage}%
  \renewcommand{\headrulewidth}{1.2pt}%
}

\fancypagestyle{style7}{%
  \fancyhf{}%
  \fancyhead[R]{نمایه}
  \fancyfoot[C]{\thepage}%
  \renewcommand{\headrulewidth}{1.2pt}%
}

\fancypagestyle{style8}{%
  \fancyhf{}%
  \fancyhead[R]{منابع و مراجع}
  \fancyfoot[C]{\thepage}%
  \renewcommand{\headrulewidth}{1.2pt}%
}
\fancypagestyle{style9}{%
  \fancyhf{}%
  \fancyhead[R]{واژه‌نامه‌ی فارسی به انگلیسی}
  \fancyfoot[C]{\thepage}%
  \renewcommand{\headrulewidth}{1.2pt}%
}
%


%دستور حذف نام لیست تصاویر و لیست جداول از فهرست مطالب
\newcommand*{\BeginNoToc}{%
  \addtocontents{toc}{%
    \edef\protect\SavedTocDepth{\protect\the\protect\value{tocdepth}}%
  }%
  \addtocontents{toc}{%
    \protect\setcounter{tocdepth}{-10}%
  }%
}
\newcommand*{\EndNoToc}{%
  \addtocontents{toc}{%
    \protect\setcounter{tocdepth}{\protect\SavedTocDepth}%
  }%
}
\newcounter{savepage}
\renewcommand{\listfigurename}{فهرست اشکال}
\renewcommand{\listtablename}{فهرست جداول}
%\renewcommand\cftsecleader{\cftdotfill{\cftdotsep}}
%%%%%%%%%%%%%%%%%%%%%%%%%%%%%
%%%%%%%%%%%%%%%%%%%%%%%%%%%%


\newcommand{\insertfig}[3]{%
\begin{figure}%
	\begin{center}%
		\ttfamily%
		\input{#1}%
		\label{#3}%
		\rmfamily%
	\end{center}%
	\caption{#2}%
\end{figure}%
}%

\begin{document}
\baselineskip=.75cm
\linespread{1.75}
%% -!TEX root = AUTthesis.tex
% در این فایل، عنوان پایان‌نامه، مشخصات خود، متن تقدیمی‌، ستایش، سپاس‌گزاری و چکیده پایان‌نامه را به فارسی، وارد کنید.
% توجه داشته باشید که جدول حاوی مشخصات پروژه/پایان‌نامه/رساله و همچنین، مشخصات داخل آن، به طور خودکار، درج می‌شود.
%%%%%%%%%%%%%%%%%%%%%%%%%%%%%%%%%%%%
% دانشکده، آموزشکده و یا پژوهشکده  خود را وارد کنید
\faculty{دانشکده ریاضی و علوم کامپیوتر}
% گرایش و گروه آموزشی خود را وارد کنید
\department{گرایش ریاضی}
% عنوان پایان‌نامه را وارد کنید
\fatitle{هوش مصنوعی - گزارش ۰۴ - پیاده سازی بازی دونفره با جستجوی تخاصمی
\\[.75 cm]
 پایان‌نامه}
% نام استاد(ان) راهنما را وارد کنید
\firstsupervisor{نام کامل استاد راهنما}
%\secondsupervisor{استاد راهنمای دوم}
% نام استاد(دان) مشاور را وارد کنید. چنانچه استاد مشاور ندارید، دستور پایین را غیرفعال کنید.
\firstadvisor{نام کامل استاد مشاور}
%\secondadvisor{استاد مشاور دوم}
% نام نویسنده را وارد کنید
\name{آترین }
% نام خانوادگی نویسنده را وارد کنید
\surname{حجت}
%%%%%%%%%%%%%%%%%%%%%%%%%%%%%%%%%%
\thesisdate{فروردین ۱۴۰۰}

% چکیده پایان‌نامه را وارد کنید
\fa-abstract{
در اين قسمت چكيده پایان نامه نوشته مي‌شو‌د‌.‌ چكيده بايد جامع و بيان‌كننده‌ خلاصه‌اي از اقدامات انجام‌شده باشد. در چكيده باید از ارجاع به مرجع و ذكر روابط رياضي، بيان تاريخچه و تعريف مسئله خودداري ‌شود. 
}


% کلمات کلیدی پایان‌نامه را وارد کنید
\keywords{هوش‌مصنوعی، هیوریستیک، پوشش راسی}



\AUTtitle
%%%%%%%%%%%%%%%%%%%%%%%%%%%%%%%%%%
\vspace*{7cm}
\thispagestyle{empty}
\begin{center}
\includegraphics[height=5cm,width=12cm]{besm}
\end{center}
% تاییدیه دفاع
\newpage
\thispagestyle{empty}
%\fontsize{18pt}{19pt}\selectfont

\section*{صفحه فرم ارزیابی و تصویب پایان نامه- فرم تأیید اعضاء كميته دفاع}

\fontsize{12pt}{14pt}\selectfont
%\renewcommand{\baselinestretch}{1.5}
\vspace*{1cm}
   در این صفحه فرم دفاع یا تایید و تصویب پایان نامه موسوم به فرم کمیته دفاع- موجود در پرونده آموزشی- را قرار دهید.
\vspace*{1cm}


\subsection*{نکات مهم:}
 
\begin{itemize}
\item
	نگارش پایان نامه/رساله باید به
	{\color{red}
		زبان فارسی
	}
	و بر اساس آخرین نسخه دستورالعمل و راهنمای تدوین پایان نامه های دانشگاه صنعتی امیرکبیر باشد.(دستورالعمل و راهنمای حاضر)
\item رنگ جلد پایان نامه/رساله چاپي كارشناسي، كارشناسي ارشد و دكترا  بايد به ترتيب مشكي، طوسي و سفيد رنگ باشد.  
\item چاپ و صحافی پایان نامه/رساله بصورت
{\color{red}
	پشت و رو(دورو)
}
بلامانع است و انجام آن توصيه مي شود. 
\end{itemize}
%%%%%%%%%%%%%%%%%%%%%%%%%%%%%%%%%%%%%%%%%%%%%%%%%%%%%%%%%%%%%%%%%%%%%%%%%%%%%%%%%%%%%%%%%%%%%%%%%%
%%%%%%%%%%%%%%%%%%%%%%%%%%%%%%%%%%%%%%%%%%%%%%%%%%%%%%%%%%%%%%%%%%%%%%%%%%%%%%%%%%%%%%%%%%%%%%%%%%
\newpage
\thispagestyle{empty}
\begin{picture}(50,50)
  \put(17,0){\includegraphics[scale=1.1]{fa-logo}}
  \put(4.5,-13){\footnotesize{دانشگاه صنعتی امیرکبیر}}
  \put(10.5,-27){\footnotesize{(پلی‌تکنیک تهران)}}
  \put(170,30){\bf{به نام خدا}}
  \put(140,-5){\Large\bf{تعهدنامه اصالت اثر}}
  \put(310,0){تاریخ: \datethesis}
\end{picture}

\vspace*{2.5cm}

اينجانب {\bf{\fname\lname}} متعهد می‌شوم که مطالب مندرج در این پایان‌نامه حاصل کار پژوهشی اینجانب تحت نظارت و راهنمایی اساتید دانشگاه صنعتی امیرکبیر بوده و به دستاوردهای دیگران که در این پژوهش از آنها استفاده شده است مطابق مقررات و روال متعارف ارجاع و در فهرست منابع و مآخذ ذکر گردیده است. این پایان‌نامه قبلاً برای احراز هیچ مدرک هم‌سطح یا بالاتر ارائه نگردیده است.

در صورت اثبات تخلف در هر زمان، مدرک تحصیلی صادر شده توسط دانشگاه از درجه اعتبار ساقط بوده و دانشگاه حق پیگیری قانونی خواهد داشت.


کلیه نتایج و حقوق حاصل از این پایان‌نامه متعلق به دانشگاه صنعتی امیرکبیر می‌باشد. هرگونه استفاده از نتایج علمی و عملی، واگذاری اطلاعات به دیگران یا چاپ و تکثیر، نسخه‌برداری، ترجمه و اقتباس از این پایان نامه بدون موافقت کتبی دانشگاه صنعتی امیرکبیر ممنوع است. 
نقل مطالب با ذکر مآخذ بلامانع است.\\
\vspace{2.5cm}


{\centerline {\bf{\fname\lname}}}
\vspace*{.2cm}
{\centerline{امضا}}
%%%%%%%%%%%%%%%%%%%%%%%%%%%%%%%%%
% چنانچه مایل به چاپ صفحات «تقدیم»، «نیایش» و «سپاس‌گزاری» در خروجی نیستید، خط‌های زیر را با گذاشتن ٪  در ابتدای آنها غیرفعال کنید.
% پایان‌نامه خود را تقدیم کنید
% نیایش خود را در فایل زیر بنویسید.
%\begin{acknowledgementpage}

\vspace{1.5cm}

{\nastaliq
{
 نويسنده پايان‌نامه، درصورت تمايل ميتواند برای سپاسگزاری پايان‌نامه خود را به شخص يا اشخاص و يا ارگان خاصی تقدیم نماید.
}}\end{acknowledgementpage}
\newpage
% سپاسگزاری را در فایل زیر بنویسید.
%%%%%%%%%%%%%%%%%%%%%%%%%%%%%%%%%%%%%
\newpage\thispagestyle{empty}
% سپاس‌گزاری
{\nastaliq
سپاس‌گزاری
}
\\[2cm]

 نويسنده پايان‌نامه می‌تواند مراتب امتنان خود را نسبت به استاد راهنما و استاد مشاور و یا ديگر افرادي كه طي انجام پايان‌نامه به نحوي او را یاری و یا با او همكاری نموده‌اند ابراز دارد.














% با استفاده از دستور زیر، امضای شما، به طور خودکار، درج می‌شود.
\signature








%%%%%%%%%%%%%%%%%%%%%%%%%%%%%%%%%%%%%%%%%
%%%%%%%%%%%%%%%%%%%%%%%%%%%%%%%%%کدهای زیر را تغییر ندهید.
\newpage\clearpage

\pagestyle{style2}

\vspace*{-1cm}
\section*{\centering چکیده}
%\addcontentsline{toc}{chapter}{چکیده}
\vspace*{.5cm}
\ffa-abstract
\vspace*{2cm}


{\noindent\large\textbf{واژه‌های کلیدی:}}\par
\vspace*{.5cm}
\fkeywords
% دستور زیر برای شماره گذاری صفحات قبل از فصل اول با حروف ابجد است.
\pagenumbering{alph}
%-----------------------------------------------------------------------------
% فایل زیر دستورات مربوط به نمایش صفحات فهرست مطالب- فهرست اشکال و جداول است.
%{\pagestyle{style2}
%\tableofcontents}\newpage
%
%\listoffigures
\cleardoublepage
\pagestyle{style6}
\tableofcontents
\pagestyle{style6}
\cleardoublepage
%اگر لیست تصاویر و لیست جداول ندارید ، کدهای زیر را با گذاشتن % در ابتدای آنها، غیرفعال کنید.
\BeginNoToc
%============
\addtocontents{lof}{\lofheading}% add heading to the first page in LoF
\pagestyle{style5}
\listoffigures
\thispagestyle{style5}
\cleardoublepage
%============
\addtocontents{lot}{\lotheading}% add heading to the first page in LoT
\thispagestyle{style4}
\listoftables
\thispagestyle{style4}
%============
%\cleardoublepage
%
\cleardoublepage
\setcounter{savepage}{\arabic{page}}
\mainmatter
\addtocontents{toc}{\tocheading}% add heading to the first page in ToC, after frontmatter entries
\EndNoToc
% در صورت تمایل می‌توانید با فعال کردن دستور بالا، لیست تصاویر را به  پایان‌نامه خود اضافه کنید.
%-------------------------------------------------------------------------symbols(فهرست نمادها)
% وجود لیست نمادها الزامیست.(لطفاً نمادهای خود را جایگذین نمادهای پیش‌فرض کنید.)
%%%%%%%%%%%%%

{\centering\LARGE\textbf{فهرست نمادها}\par}%

\pagenumbering{alph}
\setcounter{page}{\thesavepage}
%\setcounter{page}{6}
\vspace*{1cm}

\pagestyle{style3}
%\thispagestyle{empty}
%\addcontentsline{toc}{chapter}{فهرست نمادها}
\symb{\text{ نماد}}{مفهوم}
\\
%مقادیر بالا را تغییر ندهید
%%%%%%%%%%%%%%%%%%%%%%%%%%%%%%%%%%%%%%%%%%%%%%%%%%%%%%%%%
\symb{n(G)}{
تعداد رئوس گراف $G$
}
\symb{deg(v)}{
درجه‌ی راس $v$ در گراف
}
\symb{nei(v)}{
همسایه‌های راس $v$ در گراف
}
%%%%%%%%%%%%%%%%%%%%%%%%%%%%%%%%%%%%%%%

\thispagestyle{style3}
\newpage
%\pagestyle{style1}
%%%%%%%%%%%%%%%%%%%%%%%%%%%%%%%%%%%%


\pagenumbering{arabic}
\pagestyle{style1}
%--------------------------------------------------------------------------chapters(فصل ها)
\chapter{شرح مسئله و روند کار}

\section{مقدمه}

در این گزارش با استفاده از
\verb;Collaborative filtering;
یک
سیستم توصیه‌گر فیلم
می‌سازیم.
پیاده سازی و کد در
\href{https://github.com/atrin-hojjat/Uni-AI-Course-Reports/blob/main/Report\%2006/}{اینجا}
قابل مشاهده می‌باشد.
برای پیاده سازی
از
\lr{Python}
و کتاب‌خانه‌های
\lr{numpy, sklearn, scipy, pandas}
استفاده شده.

\section{روش اجرا}
برای اجرای برنامه نیاز به 
\verb;Python;
با حداقل ورژن 
$3.8$
دارید.

برای اجرای برنامه ابتدا نیاز به راه اندازی یک
\lr{Virtual environment}
برای 
\verb;Python;
است.
به این منظور در فولدر
\lr{Report 06/Codes}
دستورات زیر را اجرا کنید.

\lstset{language=bash}
\begin{latin}
\begin{lstlisting}
cd Report\\ 06/Codes
python3 -m venv venv
\end{lstlisting}
\end{latin}

برای 
\verb;activate;
کردن با توجه به سیتم عامل دستورات
\href{https://packaging.python.org/guides/installing-using-pip-and-virtual-environments/}{اینجا}
\footnote{برای جزئیات بیشتر به 
\href{https://docs.python.org/3/library/venv.html}{اینجا}
 مراجعه کنید
 }
را اجرا کنید.

برای نصب پیشنیاز ها دستورات زیر را اجرا کنید.

\begin{latin}
\begin{lstlisting}
python3 -m pip install -r requirements.txt
\end{lstlisting}
\end{latin}


برای اجرای کد نیاز به دانلود‌ دیتا‌ست دارید
\cite{10.1145/2827872}
.
به صورت پیش‌فرض برنامه درست
\lr{../Datasets/ml-latest-small/}
به دنبال دیتیست می‌گردد. برای تغییر یک
\lr{Environment variable}
با نام
\verb;DATASET_FOLDER;
درست کنید و مقدار آن را به موقعییت فولدر دیتاست تغییر دهید.
% یک فایل
% \verb;.env;
% مانند فایل
% \verb;.env.example;
% درست کنید و رد آن آدرس دانلود دیتاست را وارد کنید.

سپس با اجرای فایل
\lr{GenerateTestValues.py}
فایل‌های تست و ترینینگ ساخته می‌شوند و می‌توانید فایل
\lr{ItemToItemCF.py}
را اجرا کنید.

\begin{latin}
\begin{lstlisting}
  python3 GenerateTestValues.py
  python3 ItemToItemCF.py
\end{lstlisting}
\end{latin}

\chapter{روش تخمین امتیاز فیلم‌‌ها}
برای تخمین از تکنیک
\lr{Collaborative Filtering}
با الگوریتم
\lr{Nearest Neighbor}
برای پیداکردن فیلم‌های مشابه استفاده می‌شود.

مرجع فاصله در این گزارش زاویه بین دو بردار در نظر گرفته شده و با تا تابع
\verb;cosine;
اندازه‌گیری می‌شود. داریم

\begin{latin}
    \begin{equation}
        \cos(x, y) = \frac{x.y} {\norm{x}.\norm{y}}
    \end{equation}
\end{latin}

فرض کنید ماتریس
\verb;T;
حاوی اطلاعات رای دهی کاربران باشد و ماتریس
\verb;D;
حاوی ضرب داخلی بردار‌های امتیاز‌های هر جفت از فیلم ها باشد.
$D_{diag}$
ماتریس مربعی شامل قطر
\lr{D}
است.

\begin{latin}
    \begin{equation}
      D = T . T^{T}
    \end{equation}
\end{latin}
\begin{latin}
    \begin{equation}
    \begin{aligned}
      1 / D_{diag} &=
        \begin{pmatrix}
        1/D_{1,1} & 0 & \cdots & 0 \\
        0 & 1/D_{2,2} & \cdots & - \\
        \vdots  & \vdots  & \ddots & \vdots  \\
        0 & 0 & \cdots & 1/D_{n,n}
        \end{pmatrix} \\
        &=
        \begin{pmatrix}
        1/\norm{T_{1}} & 0 & \cdots & 0 \\
        0 & 1/\norm{T_{2}} & \cdots & - \\
        \vdots  & \vdots  & \ddots & \vdots  \\
        0 & 0 & \cdots & 1/\norm{T_{n}}
        \end{pmatrix}
    \end{aligned}
    \end{equation}
\end{latin}
\begin{latin}
    \begin{equation}
    \begin{aligned}
      X &= (1/D_{diag}).(1/D_{diag})^{T} \\
      &= \begin{pmatrix}
          \frac{1} {\norm{T_{1}}^{2}} & \frac{1}{\norm{T_{1}}\norm{T_{2}}} & \cdots & \frac{1}{\norm{T_{1}}\norm{T_{n}}} \\
          \frac{1}{\norm{T_{2}}\norm{T_{1}}} & \frac{1}{\norm{T_{2}}^{2}} & \cdots & \frac{1}{\norm{T_{2}}\norm{T_{n}}} \\
        \vdots  & \vdots  & \ddots & \vdots  \\
          \frac{1}{\norm{T_{n}}\norm{T_{1}}} & \frac{1}{\norm{T_{n}}\norm{T_{2}}} & & \cdots \frac{1}{\norm{T_{n}}^{2}} \\
        \end{pmatrix} \\
    \end{aligned}
    \end{equation}
\end{latin}
\begin{latin}
    \begin{equation}
    \begin{aligned}
      N &= \sqrt{X} = \sqrt{(1/D_{diag}).(1/D_{diag})^{T}} \\
      &= \begin{pmatrix}
          \frac{1} {\norm{T_{1}}} & \frac{1}{\sqrt{\norm{T_{1}}\norm{T_{2}}}} & \cdots & \frac{1}{\sqrt{\norm{T_{1}}\norm{T_{n}}}} \\
          \frac{1}{\sqrt{\norm{T_{2}}\norm{T_{1}}}} & \frac{1}{\norm{T_{2}}} & \cdots & \frac{1}{\sqrt{\norm{T_{2}}\norm{T_{n}}}} \\
        \vdots  & \vdots  & \ddots & \vdots  \\
          \frac{1}{\sqrt{\norm{T_{n}}\norm{T_{1}}}} & \frac{1}{\norm{T_{\sqrt{}}\norm{T_{2}}}} & & \cdots \frac{1}{\norm{T_{n}}} \\
        \end{pmatrix} \\
    \end{aligned}
    \end{equation}
\end{latin}

حال با ضرب درایه در درایه ماتریس‌های
\verb;N;
و
\verb;D;
می‌توان فاصله‌های دوبه‌دو‌ی هر دو فیلم را به‌دست آورد.

این بخش در پایتون به صورت زیر به‌دست می‌آید:

\begin{latin}
  \begin{python}
training_np = training_matrix.to_numpy()
distances = np.matmul(training_np, training_np.transpose())

distances_diagonal = 1. / np.diagonal(distances).copy()

norm_size_mul = np.sqrt(np.outer(distances_diagonal, distances_diagonal))

cosine_distances = np.multiply(distances, norm_size_mul)

distances_pd = pd.DataFrame(cosine_distances)

  \end{python}
\end{latin}

\section{پیداکردن \lr{\ttfamily{Nearest Neighbor}}}

\begin{latin}
  \begin{python}
NN = NearestNeighbors(n_neighbors=n_neighbors, metric="cosine")

NN.fit(training_csr)

...

dists, neighbors = NN.kneighbors(movie.values.reshape(1, -1), n_neighbors=n_neighbors)

  \end{python}
\end{latin}

سپس همسایه‌هایی که
\verb;user;
به آنها رای نداده را حذف می‌کنیم.

\begin{latin}
  \begin{python}
neighbors = [neighbor for neighbor in neighbors[0] if
        training_matrix.iloc[neighbor, userInd] > 0]
  \end{python}
\end{latin}

سپس باید با فرمول زیر تخمینی برای امتیاز فیلم به‌دست آوریم.

\begin{latin}
  \begin{equation}
    r_{ix}=b_{ix} + \frac{\sum_{j \in N(i; x)} D_{ij}(r_{jx} - b{jx})}{\sum_{j \in N(i; x)} D_{ij}}
  \end{equation}
\end{latin}

که در آن
$r_{ij}$
امتیاز فرد
$j$
به فیلم
$i$
و
$b_{ij}=\mu + b_{i} + b_{j}$
می‌باشد.
$\mu$
میانگین همه‌ی امتیاز‌ها،
$b_{i}$
فاصله‌ی میانگین امتیاز‌های فیلم
$i$
از میانگین کل، و
$b_{j}$
فاصله‌ی میانگین امتیاز‌های فرد
$j$
از میانگین کل است.

\begin{latin}
  \begin{python}

user_avg = training.groupby(by="userId").mean()
movie_avg = training.groupby(by="movieId").mean()
global_avg = training.mean()['rating']

bases = (np.ones(training_matrix.values.shape) * -global_avg +
        user_avg.rating.values +
        movie_avg.rating.values.reshape((movie_avg.rating.values.shape[0], 1)))

...

baseline = (global_avg +
        movie_avg.iloc[movie_avg.index==int(row['movieId'])].values[0] +
        user_avg.iloc[user_avg.index==int(row['userId'])].values[0])

neigh_dists_all = distances_pd.iloc[neighbors]
neigh_dists = neigh_dists_all.iloc[:, movieInd]


bases_nec = bases_pd.iloc[neighbors].loc[:,
        userInd]

ratings_nec = training_matrix.iloc[neighbors].iloc[:,
        training_matrix.columns==int(row['userId'])]


...


ans = bases[movieInd, userInd] + np.inner(ratings_nec.values.reshape(1, -1) -
        bases_nec.values.reshape(1, -1),
        neigh_dists.values.reshape(1, -1)) / np.sum(neigh_dists.values)


  \end{python}
\end{latin}


برای تخمین
\verb;error;
حدود
$10\%$
از دیتا را ابتدا جدا و به‌عنوان تست استفاده می‌کنیم.

\begin{latin}
  \begin{python}
import pandas as pd
import numpy as np
import os
#  from dotenv import load_dotenv

#  load_dotenv()

dataset_folder = os.getenv("DATASET_FOLDER", '../Datasets/ml-latest-small/')


def generate_tests_for_movies():
    ratings = pd.read_csv(os.path.join(dataset_folder, "ratings.csv"))

    users = ratings.groupby(by="userId").count()
    movies = ratings.groupby(by="movieId").count()

    test_users = users.sample(frac=0.1)
    test_movies = movies.sample(frac=0.1)

    test_ratings = ratings[ratings.userId.isin(test_users.index) &
            ratings.movieId.isin(test_movies.index)]
    training_ratings = ratings[(~ratings.userId.isin(test_users.index)) |
            ~ratings.movieId.isin(test_movies.index)]


    test_ratings.to_csv(os.path.join(dataset_folder, 'ratings.test.csv'))
    training_ratings.to_csv(os.path.join(dataset_folder, 'ratings.training.csv'))

if __name__ == '__main__':
    generate_tests_for_movies()

  \end{python}
\end{latin}




سپس با استفاده از این دیتا خطای مجموع مربعات را محاسبه می‌کنیم.

\begin{latin}
  \begin{equation}
    err = \sqrt{\sum_{t \in tests}{(r_{t} - expected_{t})^{2}}}
  \end{equation}
\end{latin}

\begin{latin}
  \begin{python}

    ...
    err += (ans - row['rating']) ** 2
    ...
ans = math.sqrt(err / cnt)

  \end{python}
\end{latin}

کد نهایی به صورت زیر می‌باشد:

\begin{latin}
  \begin{python}
import pandas as pd
import numpy as np
from sklearn.neighbors import NearestNeighbors
from scipy.sparse import csr_matrix
import math
import os
#  from dotenv import load_dotenv

#  load_dotenv()

dataset_folder = os.getenv("DATASET_FOLDER", '../Datasets/ml-latest-small/')

training = pd.read_csv(os.path.join(dataset_folder, "ratings.training.csv"))
test = pd.read_csv(os.path.join(dataset_folder, "ratings.test.csv"))

training_matrix = training.pivot(index="movieId", columns="userId",
        values="rating").fillna(0)
test_matrix = test.pivot(index="movieId", columns="userId",
        values="rating").fillna(0)

training_csr = csr_matrix(training_matrix)

n_neighbors = 10
NN = NearestNeighbors(n_neighbors=n_neighbors, metric="cosine")

NN.fit(training_csr)

user_avg = training.groupby(by="userId").mean()
movie_avg = training.groupby(by="movieId").mean()
global_avg = training.mean()['rating']

training_np = training_matrix.to_numpy()
distances = np.matmul(training_np, training_np.transpose())

distances_diagonal = 1. / np.diagonal(distances).copy()

norm_size_mul = np.sqrt(np.outer(distances_diagonal, distances_diagonal))

cosine_distances = np.multiply(distances, norm_size_mul)

distances_pd = pd.DataFrame(cosine_distances)

bases = (np.ones(training_matrix.values.shape) * -global_avg +
        user_avg.rating.values +
        movie_avg.rating.values.reshape((movie_avg.rating.values.shape[0], 1)))

bases_pd = pd.DataFrame(bases)

err, cnt = 0, 0
derr, dcnt = 0, 0

for index, row in test.iterrows():
    cnt += 1
    if not (training_matrix.index.values==int(row['movieId'])).any():
        print(global_avg, row['rating'])
        err += (global_avg - row['rating']) ** 2
        continue
    if not (training_matrix.index.values==int(row['userId'])).any():
        print(global_avg, row['rating'])
        err += (global_avg - row['rating']) ** 2
        continue
    movieInd = np.where(training_matrix.index.values==int(row['movieId']))[0][0]
    userInd = np.where(training_matrix.columns.values==int(row['userId']))[0][0]
    movie = training_matrix.iloc[training_matrix.index==int(row['movieId'])]

    dists, neighbors = NN.kneighbors(movie.values.reshape(1, -1), n_neighbors=n_neighbors)


    #  for neighbor in neighbors[0]:
    #      print(neighbor)

    #      print(userInd)
    #      print(training_matrix.iloc[neighbor, userInd])
    #      print(training_matrix.iloc[neighbor,
    #              training_matrix.columns==int(row['userId'])])

    neighbors = [neighbor for neighbor in neighbors[0] if
            training_matrix.iloc[neighbor, userInd] > 0]

    baseline = (global_avg +
            movie_avg.iloc[movie_avg.index==int(row['movieId'])].values[0] +
            user_avg.iloc[user_avg.index==int(row['userId'])].values[0])


    if len(neighbors) == 0:
        print(global_avg, row['rating'])
        err += (global_avg - row['rating']) ** 2
        continue


    neigh_dists_all = distances_pd.iloc[neighbors]
    neigh_dists = neigh_dists_all.iloc[:, movieInd]

    #  print("**********")
    #  print()
    #  print()

    #  print(row)

    #  print(neighbors)
    #  print(neigh_dists)

    bases_nec = bases_pd.iloc[neighbors].loc[:,
            userInd]

    ratings_nec = training_matrix.iloc[neighbors].iloc[:,
            training_matrix.columns==int(row['userId'])]


    ans = bases[movieInd, userInd] + np.inner(ratings_nec.values.reshape(1, -1) -
            bases_nec.values.reshape(1, -1),
            neigh_dists.values.reshape(1, -1)) / np.sum(neigh_dists.values)

    #  print("Base", bases[movieInd, userInd])
    #  print("Rating vec", ratings_nec.values.reshape(1, -1))
    #  print("Base vec",bases_nec.values.reshape(1, -1))
    #  print("Distances", neigh_dists.values.reshape(1, -1))
    #  print("Inner", np.inner(ratings_nec.values.reshape(1, -1) -
    #          bases_nec.values.reshape(1, -1),
    #          neigh_dists.values.reshape(1, -1)))
    #  print("Dividend", np.sum(neigh_dists.values))


    print(ans, row['rating'])
    err += (ans - row['rating']) ** 2
    derr += (ans - row['rating']) ** 2
    dcnt += 1
    #  print()
    #  print()

print("Error with mean cases: ", math.sqrt(err / cnt))
print("Error on calculated cases: ", math.sqrt(derr / dcnt))

  \end{python}
\end{latin}


\chapter{بررسی دقت با تغییر تعداد همسایه‌ها}

به‌منظور بررسی تاثیر تعداد همسایه‌ها از کد
\verb;TestNumNeighbors.py;
استفاده شده که میزان خطا را با تغییر تعداد همسایه‌ها می‌سنجد.
برای اجرای این کد با درست کردن
\lr{Environment Variable}
با نام
\verb;OUTPUT;
می‌توان خروجی
\verb;.jpg,.pgf;
از نمودار حاصل دریافت کرد.

\begin{latin}
  \begin{python}

import pandas as pd
import numpy as np
from sklearn.neighbors import NearestNeighbors
from scipy.sparse import csr_matrix
import math
import os
#  from dotenv import load_dotenv
from matplotlib.legend_handler import HandlerLine2D, HandlerTuple
import matplotlib
#  matplotlib.use("pgf")
import matplotlib.pyplot as plt
import matplotlib.lines as mlines
import seaborn as sns


#  load_dotenv()

dataset_folder = os.getenv("DATASET_FOLDER", '../Datasets/ml-latest-small/')

training = pd.read_csv(os.path.join(dataset_folder, "ratings.training.csv"))
test = pd.read_csv(os.path.join(dataset_folder, "ratings.test.csv"))

training_matrix = training.pivot(index="movieId", columns="userId",
        values="rating").fillna(0)
test_matrix = test.pivot(index="movieId", columns="userId",
        values="rating").fillna(0)

training_csr = csr_matrix(training_matrix)


user_avg = training.groupby(by="userId").mean()
movie_avg = training.groupby(by="movieId").mean()
global_avg = training.mean()['rating']

training_np = training_matrix.to_numpy()
distances = np.matmul(training_np, training_np.transpose())

distances_diagonal = 1. / np.diagonal(distances).copy()

norm_size_mul = np.sqrt(np.outer(distances_diagonal, distances_diagonal))

cosine_distances = np.multiply(distances, norm_size_mul)

distances_pd = pd.DataFrame(cosine_distances)

bases = (np.ones(training_matrix.values.shape) * -global_avg +
        user_avg.rating.values +
        movie_avg.rating.values.reshape((movie_avg.rating.values.shape[0], 1)))

bases_pd = pd.DataFrame(bases)


dataX, dataErr, dataDErr = [], [], []

for n_neighbors in range(1, 100, 3):
    err, cnt = 0, 0
    derr, dcnt = 0, 0

    NN = NearestNeighbors(n_neighbors=n_neighbors, metric="cosine")

    NN.fit(training_csr)
    for index, row in test.iterrows():
        cnt += 1
        if not (training_matrix.index.values==int(row['movieId'])).any():
            #  print(global_avg, row['rating'])
            err += (global_avg - row['rating']) ** 2
            continue
        if not (training_matrix.index.values==int(row['userId'])).any():
            #  print(global_avg, row['rating'])
            err += (global_avg - row['rating']) ** 2
            continue
        movieInd = np.where(training_matrix.index.values==int(row['movieId']))[0][0]
        userInd = np.where(training_matrix.columns.values==int(row['userId']))[0][0]
        movie = training_matrix.iloc[training_matrix.index==int(row['movieId'])]

        dists, neighbors = NN.kneighbors(movie.values.reshape(1, -1), n_neighbors=n_neighbors)


        #  for neighbor in neighbors[0]:
        #      print(neighbor)

        #      print(userInd)
        #      print(training_matrix.iloc[neighbor, userInd])
        #      print(training_matrix.iloc[neighbor,
        #              training_matrix.columns==int(row['userId'])])

        neighbors = [neighbor for neighbor in neighbors[0] if
                training_matrix.iloc[neighbor, userInd] > 0]

        baseline = (global_avg +
                movie_avg.iloc[movie_avg.index==int(row['movieId'])].values[0] +
                user_avg.iloc[user_avg.index==int(row['userId'])].values[0])


        if len(neighbors) == 0:
            #  print(global_avg, row['rating'])
            err += (global_avg - row['rating']) ** 2
            continue


        neigh_dists_all = distances_pd.iloc[neighbors]
        neigh_dists = neigh_dists_all.iloc[:, movieInd]

        #  print("**********")
        #  print()
        #  print()

        #  print(row)

        #  print(neighbors)
        #  print(neigh_dists)

        bases_nec = bases_pd.iloc[neighbors].loc[:,
                userInd]

        ratings_nec = training_matrix.iloc[neighbors].iloc[:,
                training_matrix.columns==int(row['userId'])]


        ans = bases[movieInd, userInd] + np.inner(ratings_nec.values.reshape(1, -1) -
                bases_nec.values.reshape(1, -1),
                neigh_dists.values.reshape(1, -1)) / np.sum(neigh_dists.values)

        #  print("Base", bases[movieInd, userInd])
        #  print("Rating vec", ratings_nec.values.reshape(1, -1))
        #  print("Base vec",bases_nec.values.reshape(1, -1))
        #  print("Distances", neigh_dists.values.reshape(1, -1))
        #  print("Inner", np.inner(ratings_nec.values.reshape(1, -1) -
        #          bases_nec.values.reshape(1, -1),
        #          neigh_dists.values.reshape(1, -1)))
        #  print("Dividend", np.sum(neigh_dists.values))


        #  print(ans, row['rating'])
        err += (ans - row['rating']) ** 2
        derr += (ans - row['rating']) ** 2
        dcnt += 1
        #  print()
        #  print()

    print("**************")
    print(f"Neighbors: {n_neighbors}")
    print("Error with mean cases: ", math.sqrt(err / cnt))
    print("Error on calculated cases: ", math.sqrt(derr / dcnt) if dcnt > 0 else 0)
    dataX.append(n_neighbors)
    dataErr.append(math.sqrt(err / cnt))
    dataDErr.append(False)
    dataX.append(n_neighbors)
    dataErr.append(math.sqrt(derr / dcnt) if dcnt > 0 else 0)
    dataDErr.append(True)
    print("**************")

testname = "Error by No. neighbors"
y_lab = "Mean squared error"
x_lab = "No. Neighbors"

plt.figure()
sns.relplot(data=pd.DataFrame({"n_n": dataX, "er": dataErr, "From Data": dataDErr}), hue="From Data", x="n_n", y="er", kind="line", )

plt.title(testname)
plt.xlabel(x_lab)
plt.ylabel(y_lab)

if os.getenv("OUTPUT"):
    output = os.getenv("OUTPUT")
    if not os.path.exists(os.path.dirname(os.path.join("./output",
        output, f"{testname}.jpg"))):
        try: os.makedirs(os.path.dirname(os.path.join("./output",
            output, f"{testname}.jpg")))
        except OSError as exc: # Guard against race condition
            if exc.errno != errno.EEXIST:
                raise
    plt.savefig(os.path.join("./output", output, f"{testname}.jpg"))
    matplotlib.rcParams.update({
        "pgf.texsystem": "xelatex",
        'text.usetex': True,
        'pgf.rcfonts': False,
        "font.family": "mononoki Nerd Font Mono",
        "font.serif": [],
        #  "font.cursive": ["mononoki Nerd Font", "mononoki Nerd Font Mono"],
    })
    plt.savefig(os.path.join("./output", output, f"{testname}.pgf"))

plt.show()






  \end{python}
\end{latin}


در نهایت با بررسی نمودار نهایی تفاوت چندانی بین بررسی ۱۰ تا ۱۰۰ همسایه مشاهده نشد، که البته میتواند به دلیل محدود بودن دیتا باشد.
در این نمودار
\verb;From Data;
فقط دیتا‌هایی را در نظر می‌گیرد که در همسایه‌هایش این فرد حداقل یک امتیاز داشته باشد.
در این شرایط
\verb;From Data = False;
از میانگین کلی امتیاز‌ها استفاده می‌کند.



\insertfig{figures/Error by No neighbors.pgf}{\lr{Error by No. Neighbors}}{ERR_BY_NEIGH_100}

%\include{chapter2}
%\include{chapter3}
%\include{chapter4}
%\include{chapter5}

%--------------------------------------------------------------------------appendix( مراجع و پیوست ها)
\chapterfont{\vspace*{-2em}\centering\LARGE}%

\appendix
\bibliographystyle{plain-fa}
\bibliography{references}
\chapter*{‌پیوست}
\markboth{پیوست}{}
\addcontentsline{toc}{chapter}{پیوست}
%--------------------------------------------------------------------------dictionary(واژه نامه ها)
%اگر مایل به داشتن صفحه واژه‌نامه نیستید، خط زیر را غیر فعال کنید.
%\parindent=0pt
%%
\chapter*{واژه‌نامه‌ی فارسی به انگلیسی}
\pagestyle{style9}

\addcontentsline{toc}{chapter}{واژه‌نامه‌ی فارسی به انگلیسی}
%%%%%%
\begin{multicols*}{2}

{\bf آ}
\vspace*{3mm}


\farsiTOenglish{اسکالر}{Scalar}


\vspace*{3mm}
{\bf ب}
\vspace*{3mm}

\farsiTOenglish{بالابر}{Lift}


\vspace*{3mm}
{\bf پ}
%%\vspace*{3mm}

\farsiTOenglish{پایا}{Invariant}



\vspace*{3mm}
{\bf ت}
%%\vspace*{3mm}

\farsiTOenglish{ تناظر }{Correspondence}


\vspace*{3mm}
{\bf ث}
%%\vspace*{3mm}

\farsiTOenglish{ثابت‌ساز}{Stabilizer}

\vspace*{3mm}
{\bf ج}
%%\vspace*{3mm}

\farsiTOenglish{جایگشت}{Permutation}



\vspace*{3mm}
{\bf چ}
%%\vspace*{3mm}


\farsiTOenglish{چند جمله‌ای }{Polynomial}

\vspace*{3mm}
{\bf ح}
%%\vspace*{3mm}

\farsiTOenglish{حاصل‌ضرب دکارتی}{Cartesian product}


\vspace*{3mm}
{\bf خ}
%%\vspace*{3mm}

\farsiTOenglish{خودریختی}{Automorphism}

\vspace*{3mm}
{\bf د}
%%\vspace*{3mm}

\farsiTOenglish{درجه}{Degree}


\vspace*{3mm}
{\bf ر}
%%\vspace*{3mm}


\farsiTOenglish{ریزپردازنده}{microprocessor}


\vspace*{3mm}
{\bf ز}
%%\vspace*{3mm}


\farsiTOenglish{زیرمدول}{Submodule}


\vspace*{3mm}
{\bf س}
%%\vspace*{3mm}

\farsiTOenglish{سرشت}{Character}


\vspace*{3mm}
{\bf ص}
%%\vspace*{3mm}

\farsiTOenglish{صادقانه}{Faithful}

\vspace*{3mm}
{\bf ض}
%%\vspace*{3mm}

\farsiTOenglish{ضرب داخلی}{Inner product}

\vspace*{3mm}
{\bf ط}
%%\vspace*{3mm}


\farsiTOenglish{طوقه}{Loop}


\vspace*{3mm}
{\bf ظ}
%%\vspace*{3mm}


\farsiTOenglish{ظرفیت}{Valency}
 
\vspace*{3mm}
{\bf ع}
%%\vspace*{3mm}


\farsiTOenglish{عدم مجاورت}{Nonadjacency}



\vspace*{3mm}
{\bf ف}
%%\vspace*{3mm}

\farsiTOenglish{فضای برداری}{Vector space}



\vspace*{3mm}
{\bf ک}
%%\vspace*{3mm}

\farsiTOenglish{کاملاً تحویل‌پذیر}{Complete reducibility}


\vspace*{3mm}
{\bf گ}
%%\vspace*{3mm}


\farsiTOenglish{گراف}{Graph}



\vspace*{3mm}
{\bf م}
%%\vspace*{3mm}

\farsiTOenglish{ماتریس جایگشتی}{Permutation matrix }


\vspace*{3mm}
{\bf ن}
%%\vspace*{3mm}

\farsiTOenglish{ناهمبند}{Disconnected}


\vspace*{3mm}
{\bf و}
%%\vspace*{3mm}

\farsiTOenglish{وارون‌پذیر}{Invertible}


\vspace*{3mm}
{\bf ه}
%%\vspace*{3mm}

\farsiTOenglish{همبند}{Connected}



\vspace*{3mm}
{\bf ی}
%%\vspace*{3mm}

\farsiTOenglish{یال}{Edge}




\end{multicols*}%
%%%%%%%
\chapter*{ واژه‌نامه‌ی انگلیسی به فارسی}
\pagestyle{style9}
\lhead{\thepage}\rhead{واژه‌نامه‌ی انگلیسی به فارسی}
\addcontentsline{toc}{chapter}{واژه‌نامه‌ی انگلیسی به فارسی}

\LTRmulticolcolumns
\begin{multicols}{2}
{\hfill\bf  \lr{A}}
%%\vspace*{1.5mm}

\englishTOfarsi{Automorphism}{خودریختی}

\vspace*{3mm}
{\hfill\bf   \lr{B}}
%%\vspace*{1.5mm}

\englishTOfarsi{Bijection}{دوسویی}

\vspace*{3mm}
{\hfill\bf   \lr{C}}
%%\vspace*{1.5mm}

\englishTOfarsi{Cycle group}{گروه دوری}

\vspace*{3mm}
{\hfill\bf   \lr{D}}
%%\vspace*{1.5mm}

\englishTOfarsi{Degree}{درجه}

\vspace*{3mm}
{\hfill\bf   \lr{E}}
%%\vspace*{1.5mm}

\englishTOfarsi{Edge}{یال}

\vspace*{3mm}
{\hfill\bf   \lr{F}}
%%\vspace*{1.5mm}

\englishTOfarsi{Function}{تابع}

\vspace*{3mm}
{\hfill\bf   \lr{G}}
%%\vspace*{1.5mm}

\englishTOfarsi{Group}{گروه}

\vspace*{3mm}
{\hfill\bf   \lr{H}}
%%\vspace*{1.5mm}

\englishTOfarsi{Homomorphism}{همریختی}

\vspace*{3mm}
{\hfill\bf   \lr{I}}
%%\vspace*{1.5mm}

\englishTOfarsi{Invariant}{پایا}

\vspace*{3mm}
{\hfill\bf   \lr{L}}
%%\vspace*{1.5mm}

\englishTOfarsi{Lift}{بالابر}

\vspace*{3mm}
{\hfill\bf   \lr{M}}
%%\vspace*{1.5mm}

\englishTOfarsi{Module}{مدول}

\vspace*{3mm}
{\hfill\bf   \lr{N}}
%%\vspace*{1.5mm}

\englishTOfarsi{Natural map}{نگاشت طبیعی}

\vspace*{3mm}
{\hfill\bf   \lr{O}}
%%\vspace*{1.5mm}

\englishTOfarsi{One to One}{یک به یک}

\vspace*{3mm}
{\hfill\bf   \lr{P}}
%%\vspace*{1.5mm}

\englishTOfarsi{Permutation group}{گروه جایگشتی}

\vspace*{3mm}
{\hfill\bf   \lr{Q}}
%%\vspace*{1.5mm}

\englishTOfarsi{Quotient graph}{گراف خارج‌قسمتی}

 \vspace*{3mm}
{\hfill\bf   \lr{R}}
%%\vspace*{1.5mm}

\englishTOfarsi{Reducible}{تحویل پذیر}

\vspace*{3mm}
{\hfill\bf   \lr{S}}
%%\vspace*{1.5mm}

\englishTOfarsi{Sequence}{دنباله}

 \vspace*{3mm}
{\hfill\bf   \lr{T}}
%%\vspace*{1.5mm}

\englishTOfarsi{Trivial character}{سرشت بدیهی}

\vspace*{3mm}
{\hfill\bf   \lr{U}}
%%\vspace*{1.5mm}

\englishTOfarsi{Unique}{منحصربفرد}

\vspace*{3mm}
{\hfill\bf   \lr{V}}
%%\vspace*{1.5mm}

\englishTOfarsi{Vector space}{فضای برداری}
\end{multicols}
%--------------------------------------------------------------------------index(نمایه)
%اگر مایل به داشتن صفحه نمایه نیستید، خط زیر را غیر فعال کنید.
%\pagestyle{style7}
%\printindex
%\pagestyle{style7}
%%کلمات کلیدی انگلیسی
\latinkeywords{Write a 3 to 5 KeyWords is essential. Example: AUT, M.Sc., Ph. D,..}
%چکیده انگلیسی

\en-abstract{
This page is accurate translation from Persian abstract into English.
}
%%%%%%%%%%%%%%%%%%%%% کدهای زیر را تغییر ندهید.

\newpage
\thispagestyle{empty}
\begin{latin}
\section*{\LARGE\centering Abstract}

\een-abstract

\vspace*{.5cm}
{\large\textbf{Key Words:}}\par
\vspace*{.5cm}
\elatinkeywords
\end{latin}
%% در این فایل، عنوان پایان‌نامه، مشخصات خود و چکیده پایان‌نامه را به انگلیسی، وارد کنید.
%%%%%%%%%%%%%%%%%%%%%%%%%%%%%%%%%%%%
\baselineskip=.6cm
\begin{latin}

\latinfaculty{Department of ...}


\latintitle{Title of Thesis}


\firstlatinsupervisor{Dr. }

%\secondlatinsupervisor{Second Supervisor}

\firstlatinadvisor{Dr. }

%\secondlatinadvisor{Second Advisor}

\latinname{Name}

\latinsurname{Surname}

\latinthesisdate{Month \& Year}

\latinvtitle
\end{latin}

\end{document}